

\begin{frame}{Contenu du cours}

\begin{itemize}

  \item Principes de la programmation avec BD

  \item Le langage PL/SQL
\begin{itemize}
       \item syntaxe
      \item intégration avec SQL: types, clause INTO, valeurs à NULL, parcours d'un résultat de requête
\end{itemize}
   \item Les curseurs
  \item Les triggers 
  
\end{itemize}

\vfill

Cours rédigé disponible à \url{http://sql.bdpedia.fr}{http://sql.bdpedia.fr}

\end{frame}


\begin{frame}{Rôle des procédures stockées}

\begin{itemize}

  \item SQL = langage   
  \begin{itemize}
  \item ensembliste 
  \item déclaratif (quoi, pas comment) : 
   \item pas de variables, d'itérations ou de tests, de fonctions, etc.
  \end{itemize}

  \item \textit{PL/SQL = Procedural SQL}
  
      Proposé par Oracle, existe avec des variantes dans d'autres systèmes.
      
  \item Qui : programmeur d'application, DBA
  
\end{itemize}
\end{frame}

\begin{frame}{Architectures}


\includegraphics[width=10cm]{../figures-sql/plsql}

\end{frame}

\begin{frame}{PL/SQL, quand l'utiliser, à quoi ça sert\,?}

\begin{itemize}

  \item Modes
  \begin{itemize}
  \item Interactif : on exécute du code pour, par exemple, contrôler ou corriger des données
  \item Stocké: sous forme de procédures ou de triggers
   \item Programme : appel depuis langages généralistes
  \end{itemize}

  \item But
  \begin{itemize}
  \item stocker du code "BD", dans la base
  \item factoriser certains traitements très proches des données (sécurité)
   \item performances
  \end{itemize}
\end{itemize}

\vfill

\alert{Facilite grandement l'installation de la partie applicative.}
\end{frame}



\begin{frame}[fragile]{Premier exemple, mode interactif}

\begin{footnotesize}
\begin{minted}{sql}
    DECLARE
      -- Quelques variables
      v_nbFilms   INTEGER;
      v_nbArtistes INTEGER;

    BEGIN
      -- Compte le nombre de films
      SELECT COUNT(*) INTO v_nbFilms FROM Film;
      -- Compte le nombre d'artistes
      SELECT COUNT(*) INTO v_nbArtistes FROM Artiste;

      -- Affichage des resultats
      DBMS_OUTPUT.PUT_LINE ('Nombre de films: ' || v_nbFilms);
      DBMS_OUTPUT.PUT_LINE ('Nombre d''artistes: ' || v_nbArtistes);

      EXCEPTION
        WHEN OTHERS THEN 
          DBMS_OUTPUT.PUT_LINE ('Probleme rencontre');

    END;
\end{minted}
\end{footnotesize}
\end{frame}



\begin{frame}[fragile]{Deuxième exemple, procédure stockée}

\begin{scriptsize}
\begin{minted}{sql}
    CREATE OR REPLACE PROCEDURE InsereGenre (p_genre VARCHAR) AS 
      v_genre_majuscules VARCHAR(20);
      v_count INTEGER;
      genre_existe EXCEPTION;
    BEGIN
      -- On met le parametre en majuscules
      v_genre_majuscules := UPPER(p_genre);

      -- On verifie que le genre n'existe pas 
      SELECT COUNT(*) INTO v_count 
      FROM Genre WHERE code = v_genre_majuscules;
      
      -- Si on n'a rien trouve: on insere
      IF (v_count = 0) THEN
       INSERT INTO Genre (code) VALUES (v_genre_majuscules);
      ELSE
       RAISE genre_existe;
     END IF;
  
      EXCEPTION
        WHEN genre_existe THEN
         DBMS_OUTPUT.PUT_LINE('Le genre existe deja en ' ||
                             v_count || ' exemplaire(s).');

    END;
\end{minted}
\end{scriptsize}
\end{frame}


\begin{frame}[fragile]{Troisième exemple, fonction}

\begin{scriptsize}
\begin{minted}{sql}
    CREATE OR REPLACE FUNCTION MesActeurs(v_idFilm INTEGER) 
         RETURN VARCHAR IS
      resultat VARCHAR(255);

    BEGIN
      -- Boucle prenant tous les acteurs du films 
      FOR art IN 
        (SELECT Artiste.* FROM Role, Artiste  
        WHERE idFilm = v_idFilm  AND idActeur=idArtiste)
      LOOP

        IF (resultat IS NOT NULL) THEN
          resultat := resultat || ', ' || art.prenom || ' ' || art.nom;
        ELSE   
          resultat := art.prenom || ' ' || art.nom;
       END IF;

      END LOOP;

     return resultat;
    END;
\end{minted}
\end{scriptsize}
\end{frame}

\begin{frame}[fragile]{Utilisation}

\begin{itemize}

  \item En interactif
  \begin{minted}{sql}
     SQL> start StatsFilms.sql
  \end{minted}

 \item Avec l'ordre \texttt{execute}, placé dans un autre langage (C, Java, PHP)
 
   \begin{minted}{sql}
     execute insereGenre('Policier')
  \end{minted}
 
 \item Dans une requête SQL
   \begin{minted}{sql}
      SELECT titre, MesActeurs(idFilm)  
      FROM Film WHERE idFilm=5;

TITRE        MESACTEURS(IDFILM)
------------ -----------------------------
Volte/Face   John Travolta, Nicolas Cage
  \end{minted}
 
\end{itemize}

\end{frame}

\subsection{Syntaxe de PL/SQL}



\begin{frame}[fragile]{Syntaxe PL/SQL: Les variables}

\alert{Point à retenir}\,: les types sont ceux de SQL, \alert{y compris les lignes (\texttt{RECORD})
et les tables (\texttt{TABLE}).
}
\vfill

Favorise une intégration très facile avec les requêtes.

\vfill
Exemple de déclaration d'un type \texttt{RECORD}.

\begin{scriptsize}
\begin{minted}{sql}
   TYPE adresse IS RECORD
   (no         INTEGER,
    rue        VARCHAR(40),
    ville      VARCHAR(40),
    codePostal VARCHAR(10)
   );
\end{minted}
\end{scriptsize}

On peut faire encore mieux en \alert{dérivant} un type à partir du schéma de la base.
\end{frame}


\begin{frame}[fragile]{Dérivation de types depuis le schéma}

Deux exemples pour illustrer.
\begin{itemize}

  \item \texttt{Film.titre\%TYPE} est le titre de l'attribut
     \texttt{titre} de la table \texttt{Film};

  \item \texttt{Artiste\%ROWTYPE} est  un type \texttt{RECORD} 
      correspondant aux attributs de la table \texttt{Artiste}.
\end{itemize}

\vfill

Le même principe pour les \alert{requêtes SQL} définies dans le cadre des curseurs (à suivre).

\vfill

\alert{Beaucoup plus difficile avec un langage de programmation externe.}

\end{frame}


\begin{frame}[fragile]{La section \texttt{DECLARE}}

Toutes les déclarations sont rassemblées dans la section \texttt{DECLARE}.

\begin{scriptsize}
\begin{minted}{sql}
DECLARE
  -- Constantes
  v_aujourdhui  CONSTANT DATE  := SYSDATE;
  v_pi          CONSTANT NUMBER(7,5) := 3.14116;

  -- Variables scalaires
  v_compteur    INTEGER NOT NULL := 1;
  v_nom         VARCHAR(30);

  -- Variables pour un n-uplet de la table Film et pour le resume
  v_film Film%ROWTYPE;
  v_resume Film.resume%TYPE;
\end{minted}
\end{scriptsize}

\end{frame}


\begin{frame}[fragile]{Affectation des variables}

PL/SQL propose un mécanisme d'affectation standard.

\vfill

\begin{scriptsize}
\begin{minted}{sql}
<nomVariable> := <expression>;
\end{minted}
\end{scriptsize}

\vfill

Mais surtout il est possible de placer dans des variables locales le résultat de requêtes SQL.


\begin{scriptsize}
\begin{minted}{sql}
SELECT <nomAttribut1>, [<nomAttribut2>, ...]
INTO   <nomVariable1>, [<nomVariable2>, ... ]
FROM   [...]
\end{minted}
\end{scriptsize}

\vfill

Il doit y avoir correspondance des types (facile car...)
\end{frame}


\begin{frame}[fragile]{Un exemple complet}

\begin{scriptsize}
\begin{minted}{sql}
CREATE OR REPLACE FUNCTION TitreEtMES(v_idFilm INTEGER) 
    RETURN VARCHAR IS

  -- Variables
  v_titre Film.titre%TYPE;
  v_idMES Film.idMES%TYPE;
  v_mes  Artiste%ROWTYPE;

BEGIN
  -- Recherche du film
  SELECT  titre, idMES
  INTO v_titre, v_idMES
  FROM Film
  WHERE idFilm=v_idFilm;

  -- Recherche du metteur en scene
  SELECT * INTO v_mes FROM Artiste WHERE idArtiste = v_idMES;

  return v_titre  || ', realise par ' || v_mes.prenom
                  || ' ' || v_mes.nom;
END;
\end{minted}
\end{scriptsize}

\end{frame}


\begin{frame}[fragile]{Les exceptions avec un exemple}

\begin{scriptsize}
\begin{minted}{sql}
CREATE OR REPLACE PROCEDURE TitreEnMajuscules (p_idFilm INT) AS
  v_titre Film.titre%TYPE;
  deja_fait EXCEPTION;
BEGIN
  -- Recherche du film. Exception si on ne trouve rien
  SELECT titre INTO v_titre
  FROM Film WHERE idFilm = p_idFilm;

  -- Si le titre est en majuscule, autre exception
  IF (v_titre = UPPER(v_titre)) THEN
      RAISE deja_fait;
  END IF;

  -- Mise en majuscules du titre
   UPDATE Film SET titre=UPPER(v_titre) WHERE idFilm=p_idFilm;

  EXCEPTION
    WHEN NO_DATA_FOUND THEN
     DBMS_OUTPUT.PUT_LINE('Ce film n''existe pas');
    WHEN deja_fait THEN
     DBMS_OUTPUT.PUT_LINE('Le titre est en majuscules');
    WHEN OTHERS THEN
     DBMS_OUTPUT.PUT_LINE('Autre erreur...');
END;
\end{minted}
\end{scriptsize}

\end{frame}


\begin{frame}[fragile]{Exemple d'exécution}

Illustration des levées d'exception.

\vfill

\begin{small}
\begin{minted}{sql}
SQL> execute TitreEnMajuscules(900);
Le film n'existe pas

SQL> execute TitreEnMajuscules(5);

SQL> execute TitreEnMajuscules(5);
Le titre est en majuscules
\end{minted}
\end{small}

\end{frame}

